\scanpage{56}
% archive foliation: 55

$k$ corps

\diagram{\begin{array}{lll}
\mathcal{E}_{\chi!}&\text{espaces topologiques }\inserted{\text{loc. comp.}\unsure{\text{ts}}}\text{ tels que }H^*_!(\unsure{X},k)&\text{de dim finie sur }k\\
\mathcal{E}_{\chi}&\rule{7cm}{0.4pt}\quad H^*(X,k)&\rule{4cm}{0.4pt}
\end{array}}
\margin{coh. \deleted{\unsure{faisceautique}} \inserted{singulière}}
\margin{coh. singulière}

\textcircled{a} \deleted{$\unsure{\chi(X)=\chi_!(X)}$} Si $X$ compact, $X\in\mathcal{E}_{\chi!}\Longleftrightarrow X\in\mathcal{E}_{\chi}$, et alors
\[
\chi_!(X)=\chi(X).
\]
\margin{trivial}

\textcircled{b} Si $Y$ fermé dans $X$ \inserted{loc. cpt}, alors si deux des espaces $X,Y,X-Y$ sont dans $\mathcal{E}_{\chi!}$, le troisième aussi, et on a
\[
\chi_!(X)=\chi_!(Y)+\chi_!(X-Y)
\]
(résulte de suite ex. de cohomologie à supp. cpts)

\textcircled{c} Si $X',X''$\deleted{$\unsure{X}$} fermés dans $X$ de réunion $X$, alors si trois parmi les quatre \unsure{espaces|esp.} $X,X',X'',X'\cap X''$ sont dans $\mathcal{E}_{\chi}$, le quatrième aussi, et on a
\[
\chi(X)=\chi(X')+\chi(X'')-\chi(X'\cap X'')
\]
(résulte de suite ex. de Mayer-Vietoris

\deleted{\textcircled{d}}
\[
[\text{Même}\quad X,X'\cap X''\in\mathcal{E}_{\chi}\Longrightarrow X',X''\in\mathcal{E}_{\chi}]
\]

\textcircled{d} Si $X$ \deleted{\unsure{est rétracte de}} \inserted{dominé} homotopiquement par $X'$
(\diagram{\begin{tikzcd}[column sep=large]
X \arrow[r, shift left, "i"] & X' \arrow[l, shift left, "j"]
\end{tikzcd}}
$\qquad ji\unsure{\simeq|\sim}\mathrm{id}_X\quad$) alors $X'\in\mathcal{E}_{\chi}\Longrightarrow X\in\mathcal{E}_{\chi}$, d'où
\[
(X\sim X')\unsure{\Longleftrightarrow|\Longrightarrow}(X\in\mathcal{E}_{\chi}\unsure{\Longrightarrow|\Longleftrightarrow}X'\in\mathcal{E}_{\chi}),
\]
et de plus on a alors
\[
\chi(X)=\chi(X')
\]

\textcircled{e} \deleted{\unsure{Si} $X$ \unsure{se contracte à un pt,}} $\operatorname{Card}X\leq1\Longrightarrow X\in\mathcal{E}_{\chi}$, de plus on a
\[
\chi(pt)=1
\]
(de plus $\chi(\varnothing)=0$, grâce à \unsure{e)|c)} p.ex) \margin{(trivial)}

